\subsection{Работа с событиями}

\begin{taskbox}[Реакция на событие посадки]
\indent Требуется реализовать обработчик событий, который при получении события \verb|Ev.COPTER_LANDED| полностью отключает светодиодную индикацию. Реализация не должна использовать блокирующие задержки.

\noindent\textbf{Ожидаемый результат:} после посадки все светодиоды выключены.
\end{taskbox}

\begin{taskbox}[Аварийная индикация при низком напряжении]
\indent Необходимо реализовать обработку события \verb|Ev.LOW_VOLTAGE2|, при котором все активные таймеры останавливаются, а светодиодная линейка переводится в устойчивое состояние красной индикации.

\noindent\textbf{Ожидаемый результат:} линейка постоянно светится красным цветом, сигнализируя об аварийном режиме.
\end{taskbox}

\begin{taskbox}[Индикатор загрузки (loading bar)]
\indent Требуется реализовать алгоритм постепенного заполнения светодиодной линейки от младших индексов к старшим, при котором включённые светодиоды сохраняют своё состояние.

\noindent\textbf{Ожидаемый результат:} визуальный эффект «заполняющейся» полосы.
\end{taskbox}

\begin{taskbox}[Сигнал SOS по азбуке Морзе]
\indent Необходимо реализовать световую передачу сигнала «SOS» в соответствии с азбукой Морзе: три коротких импульса, три длинных импульса и снова три коротких. Длительности импульсов должны соответствовать заданным временным параметрам.

\noindent\textbf{Ожидаемый результат:} корректно воспроизводимый световой сигнал SOS.
\end{taskbox}

\begin{taskbox}[Двоичный счётчик секунд]
\indent Требуется реализовать алгоритм отображения текущего значения секунд в двоичном виде, используя светодиоды линейки как индикаторы отдельных битов. Обновление состояния должно происходить один раз в секунду.

\noindent\textbf{Ожидаемый результат:} на линейке отображается двоичное представление времени в секундах.
\end{taskbox}
